\documentclass[12pt,a4paper]{article}
\usepackage[utf8]{inputenc}
\usepackage[spanish,es-noquoting]{babel}
\usepackage[T1]{fontenc}
\usepackage{lmodern}
\usepackage{geometry}
\geometry{top=2cm, bottom=2cm, left=2.5cm, right=2.5cm}
\usepackage{xcolor}
\usepackage{listings}
\usepackage{enumitem}
\usepackage{fancyhdr}
\usepackage{hyperref}
\usepackage{tikz}
\usetikzlibrary{arrows.meta, positioning, calc, shapes.geometric}

\definecolor{codegreen}{rgb}{0,0.6,0}
\definecolor{backcolour}{rgb}{0.95,0.95,0.92}

\lstset{
    language=C,
    backgroundcolor=\color{backcolour},
    commentstyle=\color{codegreen},
    keywordstyle=\color{blue},
    basicstyle=\ttfamily\small,
    breaklines=true,
    frame=single,
    numbers=left,
    tabsize=4,
    extendedchars=true,
    showstringspaces=false,
    inputencoding=utf8,
    literate={á}{{\'a}}1 {é}{{\'e}}1 {í}{{\'i}}1 {ó}{{\'o}}1 {ú}{{\'u}}1 {ñ}{{\~n}}1
}

\pagestyle{fancy}
\fancyhf{}
\lhead{Monitorías Sistemas Operativos 2026-I}
\rhead{Capítulo 1: Taller Práctico}
\cfoot{\thepage}

\title{\textbf{Taller Práctico: División de Tareas y Comunicación vía Archivos}}
\author{Monitorías de Sistemas Operativos}
\date{\today}

\begin{document}

\maketitle

\section*{Introducción}
En este taller, consolidaremos el conocimiento sobre la creación de procesos mediante \texttt{fork()}. El objetivo es implementar un programa que divida el procesamiento de un archivo de datos entre múltiples hijos, utilizando archivos externos como mecanismo de persistencia y comunicación.

\section*{Descripción del Problema: Auditoría de Almacenes}
La multinacional IBM acaba de finalizar la auditoría anual de sus infraestructuras logísticas a nivel mundial. Tras recolectar los registros de cada sede, han consolidado el valor total de existencias en un archivo maestro denominado \texttt{data.in}. Este archivo presenta un volumen de datos crítico: la primera línea indica el número total de registros $n$, seguida de $n$ valores enteros que representan el capital de cada almacén individual.

\vspace{0.5cm}

El departamento de ingeniería ha determinado que procesar esta información de forma puramente secuencial es inaceptable, ya que los tiempos de espera y el consumo monohilo de recursos retrasan la toma de decisiones corporativas. Por ello, se le ha otorgado acceso a las estaciones de trabajo de alto rendimiento con una consigna clara: paralelizar.

\vspace{0.5cm}

Como junior developer, su misión es diseñar una herramienta que aproveche la arquitectura del sistema mediante el uso de la llamada al sistema \texttt{fork()}. El objetivo es generar procesos asíncronos que compartan la carga de trabajo, sumen las secciones asignadas del inventario y retornen un resultado consolidado de forma ultraeficiente. El destino de la auditoría financiera está en sus manos; asegúrese de que la coordinación de estos procesos sea perfecta.

\section*{Esquema del Laboratorio}
\begin{center}
\begin{tikzpicture}[
    node distance=1.8cm and 2.5cm,
    process/.style={rectangle, draw, thick, rounded corners, fill=blue!5, minimum width=3.5cm, minimum height=1.2cm, align=center, font=\small\sffamily},
    file/.style={cylinder, draw, thick, fill=orange!10, shape border rotate=90, minimum width=2cm, minimum height=1.5cm, aspect=0.4, font=\small\sffamily},
    arrow/.style={-Stealth, thick, draw=black!70}
]
    % Nodes
    \node[process] (Padre) {Proceso Padre};
    \node[process] (H1) [below left=of Padre] {Proceso Hijo 1 \\ \textit{Suma Parte A}};
    \node[process] (H2) [below right=of Padre] {Proceso Hijo 2 \\ \textit{Suma Parte B}};
    \node[file] (Input) [right=of Padre, xshift=1cm] {data.in};
    \node[file] (Output) [below=of Padre, yshift=-1.5cm] {out.txt};

    % Connections
    \draw[arrow] (Input) -- node[above, font=\scriptsize] {Lectura inicial} (Padre);
    \draw[arrow] (Padre) -- node[left, font=\scriptsize, xshift=-0.2cm] {fork()} (H1);
    \draw[arrow] (Padre) -- node[right, font=\scriptsize, xshift=0.2cm] {fork()} (H2);
    \draw[arrow] (H1) |- node[below left, pos=0.25, font=\scriptsize] {fprintf()} (Output);
    \draw[arrow] (H2) |- node[below right, pos=0.25, font=\scriptsize] {fprintf()} (Output);
    \draw[arrow] (Output.north) -- node[right, font=\scriptsize] {Consolidación total} (Padre.south);
\end{tikzpicture}
\end{center}

\subsection*{Responsabilidades de los Procesos}
Para que el flujo sea exitoso, cada proceso debe cumplir un rol específico:


\subsection*{El proceso padre deberá:}
\begin{enumerate}
    \item Validar que el nombre del archivo de entrada sea pasado \textbf{obligatoriamente} como argumento por la línea de comandos (\texttt{argc} y \texttt{argv}).
    \item Leer en un vector los datos del archivo de entrada indicado.
    \item Estimar los índices (principio y fin) del vector sobre los cuales cada proceso hijo trabajará (usando un \textit{delta} o \textit{chunk}).
    \item Crear los procesos hijos mediante \texttt{fork()}.
    \item Esperar la terminación de los hijos mediante \texttt{wait()}.
    \item Imprimir el resultado final consolidado.
\end{enumerate}

\subsection*{Los procesos hijos deberán:}
\begin{enumerate}
    \item Recorrer el vector en las posiciones indicadas por el padre.
    \item Estimar la suma de las posiciones recorridas.
    \item Escribir el resultado parcial en el archivo de salida \texttt{out.txt}.
\end{enumerate}
\newpage

\section*{Gestión Crítica de Procesos}
Para garantizar la robustez del sistema, debe prestar especial atención a:
\begin{enumerate}
    \item \textbf{Identificación:} Use condicionales para asegurar que los hijos solo ejecuten su tarea asignada y finalicen sin continuar la ejecución del código destinado al padre.
    \item \textbf{Zombies y Huérfanos:} 
    \begin{itemize}
        \item El padre debe invocar \texttt{wait()} para recolectar el estado de finalización de sus hijos. No hacerlo generará procesos en estado \texttt{defunct} (zombies).
    \end{itemize}
\end{enumerate}

\section*{Funciones de Apoyo}
Se proporcionan las siguientes funciones para integrarlas en su desarrollo.

\subsection*{Utilidad de Error}
\begin{lstlisting}
void error(char *msg) {
    perror(msg);
    exit(-1);
}
\end{lstlisting}

\subsection*{Lectura de Datos}
\begin{lstlisting}
int read_numbers(const char *filename, int **vec) {
    int c, number, totalNumbers;
    FILE *infile = fopen(filename, "r");
    if(!infile) error("Error opening input file");
    
    fscanf(infile, "%d", &totalNumbers);
    *vec = (int *)calloc(totalNumbers, sizeof(int));
    if(!*vec) error("Error allocating memory");
    
    for(c = 0; c < totalNumbers; c++){
        if(fscanf(infile, "%d", &number) != EOF){
            (*vec)[c] = number;
            printf("%d\n", number); // Debug purposes
        }
    }
    fclose(infile);
    return c;
}
\end{lstlisting}

\newpage

\subsection*{Cálculo de Suma Final}
\begin{lstlisting}
int read_total() {
    FILE *infile;
    int sum_1 = 0, sum_2 = 0, total = 0;
    infile = fopen("out.txt", "r");
    if (!infile) error("Error opening output file for reading");
    
    fscanf(infile, "%d", &sum_1);
    fscanf(infile, "%d", &sum_2);

    total = sum_1 + sum_2;
    
    fclose(infile);
    return total;
}
\end{lstlisting}

\section*{Consideraciones Importantes}

\subsection*{Comunicación y Paso por Referencia}
La función \texttt{read\_numbers} recibe un puntero doble (\texttt{int **vec}) para permitir que los datos leídos sean accedidos por fuera de la función mediante \textit{paso por referencia}. Esto permite que la función reserve la memoria necesaria y el cambio persista en el \texttt{main}. 

Debe invocarse haciendo uso del parámetro pasado por consola:
\begin{lstlisting}
int *vector;
int total;

// Suponiendo que el archivo es el primer argumento
total = read_numbers(argv[1], &vector);
\end{lstlisting}

\subsection*{Limpieza del Entorno}
Debido a que el esquema de comunicación está basado en archivos compartidos, es necesario que \textbf{antes de cada ejecución} se elimine el archivo \texttt{out.txt} para evitar leer valores de ejecuciones previas. Para esto, utilice la llamada al sistema \texttt{remove}:
\begin{lstlisting}
remove("out.txt");
\end{lstlisting}

\end{document}
